\documentclass[11pt]{article}
\usepackage[margin=1in]{geometry}
\usepackage{amsmath,amssymb,amsthm}
\usepackage[T1]{fontenc}
\usepackage{lmodern}
\usepackage{hyperref}

\title{Literature-Implied Negative Answer for arXiv:1912.12365}
\author{fa\_banach\_001, agent\_lane\_01}
\date{2026-07-03}

\begin{document}
\maketitle

\begin{abstract}
Burton and Juschenko's arXiv:1912.12365 formulates two strengthening
conjectures around positive definite extensions and proves that they imply
Connes' embedding conjecture. Later literature, via the theorem
\( \mathrm{MIP}^*=\mathrm{RE} \), refutes Connes' embedding. Therefore the
half finite approximation conjecture and the low energy extension conjecture
of arXiv:1912.12365 are false by contrapositive.
\end{abstract}

\section*{Source conjectures}

In arXiv:1912.12365, Burton and Juschenko state Connes' embedding conjecture
as Conjecture 1.1, source label \texttt{thm.CEC}:
\[
 C^*(\mathbb F)\otimes_{\max} C^*(\mathbb F)
 =
 C^*(\mathbb F)\otimes_{\min} C^*(\mathbb F).
\]
They then state Conjecture 1.3, source label \texttt{thm.half}, on existence
of half finite approximations for representations of
\(\mathbb F\times \mathbb F\). They also state Conjecture 1.5, source label
\texttt{lem.nogain}, on extending finite families of normalized strictly
positive definite functions without increasing the pairwise energy quantities.

\section*{Implication supplied by the source paper}

The source paper explicitly supplies the implication chain needed here. Its
introduction says that the main topic is a theorem showing
\(\texttt{lem.nogain}\Rightarrow\texttt{thm.CEC}\). In the subsection
``Deducing Connes' embedding conjecture from half finite approximation
conjecture'', the authors prove
\[
 \texttt{thm.half} \Longrightarrow \texttt{thm.CEC}.
\]
The proof ends by saying that, in order to prove Connes' embedding, it
suffices to prove \texttt{thm.half}. Later, the section titled
``Proof of Conjecture \texttt{thm.half} from Conjecture
\texttt{lem.nogain}'' states that the paper proves
\[
 \texttt{lem.nogain} \Longrightarrow \texttt{thm.half}.
\]
Thus, as stated in arXiv:1912.12365,
\[
 \texttt{lem.nogain} \Longrightarrow \texttt{thm.half}
 \Longrightarrow \texttt{thm.CEC}.
\]

\section*{Later literature}

Ji, Natarajan, Vidick, Wright, and Yuen prove
\(\mathrm{MIP}^*=\mathrm{RE}\) in arXiv:2001.04383 and state that their result
refutes Connes' embedding conjecture. Goldbring's survey arXiv:2109.12682
records the negative solution to the Connes embedding problem and explains
its relationship to the Kirchberg/QWEP/Tsirelson equivalence web.

\section*{Conclusion}

The accepted later literature gives
\[
 \neg\,\texttt{thm.CEC}.
\]
Since arXiv:1912.12365 proves
\(\texttt{thm.half}\Rightarrow\texttt{thm.CEC}\), the half finite
approximation conjecture \texttt{thm.half} is false. Since the same source
proves \(\texttt{lem.nogain}\Rightarrow\texttt{thm.half}\), the low energy
extension conjecture \texttt{lem.nogain} is also false.

This is a literature-implied negative answer. It does not construct an
explicit finite counterexample to \texttt{lem.nogain}, and it does not affect
the weak extension theorem proved in the source paper.

\begin{thebibliography}{9}

\bibitem{BJ}
Peter Burton and Kate Juschenko.
\newblock \emph{Extension of positive definite functions and Connes' embedding
conjecture}.
\newblock arXiv:1912.12365.

\bibitem{MIPRE}
Zhengfeng Ji, Anand Natarajan, Thomas Vidick, John Wright, and Henry Yuen.
\newblock \emph{\(\mathrm{MIP}^*=\mathrm{RE}\)}.
\newblock arXiv:2001.04383.

\bibitem{Goldbring}
Isaac Goldbring.
\newblock \emph{The Connes Embedding Problem: A guided tour}.
\newblock arXiv:2109.12682.

\end{thebibliography}

\end{document}
